\documentclass[11pt,a4paper]{article}

% ── Packages ──────────────────────────────────────────────────────────────
\usepackage[margin=1in]{geometry}
\usepackage[utf8]{inputenc}
\usepackage[T1]{fontenc}
\usepackage{lmodern}
\usepackage{microtype}
\usepackage{amsmath,amssymb,amsthm}
\usepackage{booktabs}
\usepackage{graphicx}
\usepackage{xcolor}
\usepackage{hyperref}
\usepackage{cleveref}
\usepackage{enumitem}
\usepackage{titlesec}
\usepackage{abstract}
\usepackage{natbib}
\usepackage{tabularx}
\usepackage{array}

% ── Hyperref setup ────────────────────────────────────────────────────────
\hypersetup{
    colorlinks=true,
    linkcolor=blue!60!black,
    citecolor=blue!60!black,
    urlcolor=blue!60!black,
    pdftitle={Jasper: Engineering a Global Workspace for Latent Reasoning in Hybrid State-Space Models},
    pdfauthor={Jasper Research}
}

% ── Title ─────────────────────────────────────────────────────────────────
\title{\textbf{Jasper: Engineering a Global Workspace for Latent Reasoning\\
in Hybrid State-Space Models}}
\author{Jasper Research Project}
\date{August 2026}

% ── Custom commands ───────────────────────────────────────────────────────
\newcommand{\R}{\mathbb{R}}
\newcommand{\workspace}{\textsc{Workspace}}
\newcommand{\jspace}{\textsc{J-space}}
\newcommand{\trm}{\textsc{TRM}}

\begin{document}
\maketitle

\begin{abstract}
\noindent
Recent work from Anthropic has revealed that language models develop a compact, privileged internal workspace---the \jspace{}---that causally carries multi-step reasoning, while Samsung's Tiny Recursive Model (TRM) has demonstrated that iterative refinement through a tiny recurrent network can substitute for parameter count on hard reasoning tasks. We present Jasper, an architecture that combines these two insights into a single system: a hybrid state-space/attention backbone with an explicitly engineered workspace module and a recurrent core that iterates on the workspace state. The key architectural observation is that the SSM's fixed-size compressed state is functionally shaped like the \jspace{}---a coincidence that makes the combination more powerful than applying the same ideas to a pure transformer. The implementation uses a decayed linear attention (retention) layer in place of the original selective scan, preserving the key property---a fixed-size, input-dependent state---while being simpler and faster on the target hardware. We argue that this approach is a direct and more efficient alternative to chain-of-thought ``thinking tokens'': it reasons in latent space without generating tokens, keeps memory flat in both sequence length and loop count, and allows trading inference-time compute depth for parameter count. We describe a proof-of-concept experiment that trains three parameter-matched model variants on synthetic reasoning tasks with controllable depth, and outline the pre-registered criteria for a go/no-go decision on scaling.
\end{abstract}

\tableofcontents
\newpage

% ══════════════════════════════════════════════════════════════════════════
\section{Introduction}
% ══════════════════════════════════════════════════════════════════════════

Large language models have demonstrated impressive reasoning capabilities, but the dominant paradigm for scaling reasoning---generating long chains of ``thinking tokens'' before producing an answer---has fundamental inefficiencies. Each thinking token requires a full forward pass through the model, the key-value cache grows linearly with the length of the reasoning trace, and the entire chain must be serialized in token space, consuming both memory and latency. On constrained deployment targets such as web browsers or mobile devices, this memory growth is the binding constraint.

Two recent findings suggest an alternative path. First, Anthropic's discovery of the \jspace{} \citep{anthropic2026jspace} revealed that the reasoning-critical machinery inside a language model is already small: a compact set of verbalizable representations, occupying less than 10\% of activation variance, carries the causal bulk of multi-step reasoning. This workspace is not the full residual stream---it is a privileged, broadcasted subset that functions analogously to the global workspace in cognitive science \citep{baars2005global}. Second, Samsung's Tiny Recursive Model (TRM) \citep{jolicoeur2025trm} and Geiping et al.'s recurrent-depth transformers \citep{geiping2025recurrent} have shown that iterative refinement through a small recurrent network can substitute for parameter count: a 7M-parameter model recursing on itself can outperform models with $1000\times$ more parameters on reasoning benchmarks.

Jasper combines these insights. We engineer an explicit workspace module---a compact set of learned slot vectors that serve as a persistent scratchpad---and pair it with a recurrent core that iterates on the workspace state. The backbone is a hybrid state-space/attention model. The original design used Mamba-2's selective state-space duality (SSD) \citep{dao2024mamba2}; the current implementation uses a decayed linear attention (retention) layer that preserves the key property---a fixed-size, input-dependent compressed state---while being simpler and faster on the target hardware (Tenstorrent Blackhole). The SSM state is functionally shaped like the \jspace{}: the architecture already has a mechanism that resembles the global workspace, so the engineered workspace and the SSM state reinforce each other rather than competing for the same functional role.

The result is a model that reasons in latent space without generating tokens, keeps memory flat in both sequence length and loop count, and can trade inference-time compute (more loop iterations) for parameter count. We argue that this is a direct alternative to thinking tokens---one that is more memory-efficient, faster, and potentially more capable of reasoning that is not easily expressed in words.

% ══════════════════════════════════════════════════════════════════════════
\section{Background}
% ══════════════════════════════════════════════════════════════════════════

% ──────────────────────────────────────────────────────────────────────────
\subsection{The J-Space: A Global Workspace in Language Models}

\citet{anthropic2026jspace} introduced the Jacobian lens (J-lens), an interpretability technique that identifies which concepts a language model is \emph{poised to verbalize} at any point in its processing. Applying the J-lens across layers reveals a set of representations---the \jspace{}---that exhibits the functional properties of a global workspace \citep{dehaene2014conscious,baars2005global}:

\begin{itemize}[nosep]
    \item \textbf{Verbal report:} J-space contents can be read out as words the model could say.
    \item \textbf{Directed modulation:} Instructing the model to think about a concept causes that concept to appear in the J-space.
    \item \textbf{Silent reasoning:} Intermediate steps of multi-step problems appear in the J-space before the model verbalizes them.
    \item \textbf{Selective causality:} Ablating the J-space destroys multi-step reasoning while leaving single-hop recall intact.
    \item \textbf{Compactness:} The J-space carries fewer than $\sim$25 concepts at a time and occupies $<$10\% of activation variance.
\end{itemize}

Crucially, the \jspace{} was not designed---it emerged during training. But its properties are exactly what one would want from an engineered reasoning workspace: it is small, privileged, broadcasted, and causally responsible for multi-step reasoning. This raises the question: \emph{can we engineer such a workspace explicitly, and does doing so improve reasoning in a small model?}

% ──────────────────────────────────────────────────────────────────────────
\subsection{Iterative Reasoning with Tiny Networks}

\citet{jolicoeur2025trm} introduced the Tiny Recursive Model (TRM), a 7M-parameter network that recurses on itself to iteratively refine answers. TRM alternates between a ``think'' step (updating a latent scratchpad $\mathbf{z}$) and an ``act'' step (updating a solution embedding $\mathbf{y}$), unrolled for up to 16 iterations with deep supervision. Despite having less than 0.01\% of the parameters of frontier LLMs, TRM achieves 45\% accuracy on ARC-AGI-1 and 8\% on ARC-AGI-2, outperforming models like DeepSeek R1 and o3-mini on these benchmarks.

Independently, \citet{geiping2025recurrent} demonstrated that recurrent-depth transformers---which iterate a transformer block for a variable number of steps at inference time---can scale test-time compute in latent space. Their 3.5B-parameter model, trained on 800B tokens, improves performance on reasoning benchmarks with additional inference-time iterations, sometimes matching the performance of a 50B-parameter model.

Both works share a key insight: \emph{depth (iterations) can substitute for width (parameters)}. But both use attention-based cores, which means the key-value cache grows with sequence length, and the recurrent state is carried in the residual stream---a high-dimensional, unstructured representation.

% ──────────────────────────────────────────────────────────────────────────
\subsection{Mamba-2 and State Space Duality}

\citet{dao2024mamba2} established the state space duality (SSD) framework, showing that Mamba's selective state-space model is closely related to attention through a shared mathematical structure. The Mamba-2 layer computes:

\begin{equation}
    y_t = \sum_{s \leq t} \underbrace{\exp\!\Big(\!\!-\!\sum_{k=s+1}^{t} a_k\Big)}_{\text{decay}} \cdot \big(C_t^\top B_s\big) \cdot V_s
\end{equation}

where $B_s, C_t$ are input-dependent projections, $a_k$ is a scalar decay, and $V_s$ is the value vector. This can be viewed as a form of attention with a learned, position-dependent decay mask---the ``selective scan.''

The key property for our purposes is that the SSM maintains a \textbf{fixed-size state} $\mathbf{h}_t$ that is updated at each position:

\begin{equation}
    \mathbf{h}_t = \bar{A}_t \mathbf{h}_{t-1} + \bar{B}_t x_t
\end{equation}

where $\bar{A}_t, \bar{B}_t$ are discretized state transition matrices. This state is \textbf{compressed}: it has a fixed dimension $d_{\text{state}}$ regardless of sequence length, and it is \textbf{selective}: the input-dependent gates control what information is written and what is forgotten. This is functionally analogous to the \jspace{}---a small, privileged state that carries task-relevant information forward through the computation.

\paragraph{From Mamba-2 to retention (Mamba-3).}
The current implementation replaces the selective scan with a decayed linear attention (retention) layer, which we refer to as ``Mamba-3'' in the codebase. The retention layer computes:
\begin{equation}
    y_t = \sum_{s \leq t} \gamma^{t-s} \cdot (C_t^\top B_s) \cdot V_s
\end{equation}
where $\gamma \in (0, 1)$ is a learned per-head decay scalar. This preserves the key properties of the SSM state---fixed-size, input-dependent, compressed---while replacing the input-dependent decay $a_k$ with a simpler learned scalar $\gamma$. The decay matrix $D[t,s] = \gamma^{t-s}$ is computed on-device. This change was motivated by hardware constraints: the selective scan's causal conv1d is not well-supported on ttnn, while the retention formulation is pure matmuls. The backward pass (gradient of the loss w.r.t.\ $\gamma$) is computed via a row-sum/col-sum reduction of the $T \times T$ decay matrix, also on-device.

% ══════════════════════════════════════════════════════════════════════════
\section{The SSM Coincidence}
\label{sec:coincidence}
% ══════════════════════════════════════════════════════════════════════════

The central architectural observation behind Jasper is what we call the \emph{SSM coincidence}: the state-space model's compressed state already maintains a representation with the functional shape of the \jspace{}, without any engineering. The observation was originally made for Mamba-2's selective scan; the current retention-based implementation preserves it. Consider the parallels:

\begin{table}[h]
\centering
\small
\begin{tabularx}{\textwidth}{p{0.22\textwidth} p{0.36\textwidth} p{0.36\textwidth}}
\toprule
\textbf{Property} & \textbf{J-space (Anthropic)} & \textbf{SSM State (Mamba-2 / Retention)} \\
\midrule
Fixed size & $\sim$25 concepts, $<$10\% of activation variance & $d_{\text{state}}=64$ per head, regardless of sequence length \\
Selective & Broadcasted more widely than other representations & Input-dependent gates control write/forget \\
Persistent & Carries intermediate reasoning steps across layers & Carries information across the sequence \\
Causally responsible & Ablation destroys multi-step reasoning & Removing state connections breaks long-range dependencies \\
Emerges in training & Not designed, emerged in Claude & Not designed for reasoning, emerged from language modeling \\
\bottomrule
\end{tabularx}
\caption{Functional parallel between the \jspace{} and the SSM state (Mamba-2 selective scan or retention).}
\label{tab:coincidence}
\end{table}

This parallel is not merely metaphorical. When we build a hybrid model with SSM layers and an explicit workspace module, the SSM's compressed state and the workspace slots \emph{reinforce each other}: the SSM state provides a natural substrate for the workspace to read from and write to, and the workspace provides a structured, interpretable interface to the SSM's otherwise opaque state. In a pure transformer, the workspace must carve out its role from the residual stream against the competing pressure of attention's global mixing. In a hybrid SSM model, the workspace aligns with an existing architectural feature.

The coincidence becomes powerful when combined with a recurrent core. In a recurrent transformer \citep{geiping2025recurrent}, each iteration processes the full sequence with attention, growing the KV cache. In a recurrent SSM model, each iteration updates the SSM state in-place---memory is flat in both sequence length \emph{and} loop count. The workspace state, carried between iterations, provides a stable anchor for iterative refinement without any memory growth.

% ══════════════════════════════════════════════════════════════════════════
\section{Architecture}
\label{sec:architecture}
% ══════════════════════════════════════════════════════════════════════════

% ──────────────────────────────────────────────────────────────────────────
\subsection{Overview}

The Jasper architecture combines four components:

\begin{enumerate}[nosep]
    \item \textbf{Retention (Mamba-3) layers} for sequence processing with flat memory. The retention layer uses decayed linear attention with a learned per-head decay scalar $\gamma$, preserving the SSM's fixed-size compressed state while being pure matmuls (no selective scan or causal conv1d).
    \item \textbf{Attention layers} at selected positions for exact retrieval capability.
    \item \textbf{Workspace module} --- a perceiver-style cross-attention block with learned slot vectors.
    \item \textbf{Recurrent core} --- a range of layers that are iterated $K$ times, with the workspace state carried between iterations.
\end{enumerate}

% ──────────────────────────────────────────────────────────────────────────
\subsection{Workspace Module}
\label{sec:workspace}

The workspace module is an explicit implementation of the \jspace{} concept. It maintains $m=16$ learned slot vectors $\mathbf{S} \in \R^{m \times d}$, where $d=384$ is the model dimension. At each application, it performs two phases:

\paragraph{Read phase (sequence $\to$ workspace).}
Slots attend over the hidden states $\mathbf{X} \in \R^{T \times d}$:
\begin{align}
    \mathbf{Q}_s &= \mathbf{S}\, W_Q, \quad \mathbf{K}_x = \mathbf{X}\, W_K, \quad \mathbf{V}_x = \mathbf{X}\, W_V \\
    \mathbf{A}_{\text{read}} &= \text{softmax}\!\left(\frac{\mathbf{Q}_s \mathbf{K}_x^\top}{\sqrt{d_h}}\right) \\
    \mathbf{S}' &= \text{SlotNorm}\!\left(\mathbf{S} + \sigma(g_r) \cdot \mathbf{A}_{\text{read}} \mathbf{V}_x\, W_O\right)
\end{align}

where $d_h = d / n_{\text{heads}}$ is the per-head dimension, $\sigma$ is the sigmoid function, and $g_r$ is a learned scalar gate. The sigmoid gating allows the model to learn how much of the sequence to absorb into the workspace at each step.

\paragraph{Write phase (workspace $\to$ sequence).}
Hidden states attend over the updated slots:
\begin{align}
    \mathbf{Q}_x &= \mathbf{X}\, W_Q', \quad \mathbf{K}_s = \mathbf{S}'\, W_K', \quad \mathbf{V}_s = \mathbf{S}'\, W_V' \\
    \mathbf{A}_{\text{write}} &= \text{softmax}\!\left(\frac{\mathbf{Q}_x \mathbf{K}_s^\top}{\sqrt{d_h}}\right) \\
    \mathbf{X}' &= \text{RMSNorm}\!\left(\mathbf{X} + \sigma(g_w) \cdot \mathbf{A}_{\text{write}} \mathbf{V}_s\, W_O'\right)
\end{align}

The slot state $\mathbf{S}'$ is persistent: it is carried forward to the next workspace application and, in the recurrent core, across loop iterations. SlotNorm (RMSNorm applied to the slot dimension) prevents unbounded growth of slot activations across iterations.

\paragraph{Slot parameter normalization.}
The learned slot parameters $\mathbf{S}$ are subject to a positive feedback loop: large slots produce sharp attention distributions, which produce large gradients, which cause the slots to grow further. To break this loop, the slot parameters are normalized to unit RMS after each optimizer step:
\begin{equation}
    \mathbf{S} \leftarrow \frac{\mathbf{S}}{\text{RMS}(\mathbf{S})}, \quad \text{RMS}(\mathbf{S}) = \sqrt{\frac{1}{md}\sum_{i,j} S_{ij}^2}
\end{equation}
This is a parameter constraint (analogous to weight clipping in GANs) rather than a learned normalization. It allows the \emph{direction} of each slot vector to change freely while constraining its magnitude, preventing the runaway growth that causes gradient explosion.

\paragraph{Spectral normalization of workspace weights.}
Even with slot normalization, the attention weight matrices ($W_Q, W_K, W_V, W_O$) can grow during training, producing large attention logits that create sharp softmax distributions. To address this root cause directly, we cap the spectral norm of all eight workspace weight matrices at a bound $C$ after each optimizer step:
\begin{equation}
    W \leftarrow W \cdot \min\!\left(1,\; \frac{C}{\sigma_{\max}(W)}\right)
\end{equation}
where $\sigma_{\max}(W)$ is the largest singular value, estimated via 10 steps of power iteration on-device. This is the technique that stabilized GAN discriminators \cite{miyato2018spectral}, which face the same problem: unbounded weights $\to$ sharp distributions $\to$ gradient instability. The cap (rather than a hard normalization to $C$) lets the optimizer grow weights naturally up to the bound, then prevents further growth — breaking the positive feedback loop without disrupting training.

With $C = 5$, unit-RMS slots, and LayerNorm on the input, the attention logits are bounded by:
\begin{equation}
    |\text{logit}_{ij}| \leq \|W_Q\|_2 \|W_K\|_2 \|\text{slots}\| \|x\| / \sqrt{d_h} \leq C^2 / \sqrt{d_h} \approx 2.55
\end{equation}
giving a maximum attention ratio of $e^{2 \times 2.55} \approx 164{:}1$ — sufficient for selective attention, but preventing the $e^{100+}$ ratios that cause gradient explosion. The model can still learn the \emph{direction} of each weight matrix (which determines what it attends to) and grow its magnitude up to $C$ (which determines attention sharpness), but cannot grow beyond $C$.

\paragraph{Design rationale.}
The workspace is deliberately small ($m=16$ slots, $\sim$2M parameters total) to mirror the compactness of the \jspace{}. The two-phase design (read then write) mirrors the global workspace theory's broadcast cycle: information is first consolidated into the workspace, then broadcast back to the processing stream.

% ──────────────────────────────────────────────────────────────────────────
\subsection{Recurrent Core}
\label{sec:recurrent-core}

Layers $c_{\text{start}}$ through $c_{\text{end}}-1$ (by default, layers 6--9) form the recurrent core. These layers are applied $K$ times in sequence, with the workspace state carried between iterations:

\begin{equation}
    \mathbf{x}^{(k+1)}, \mathbf{S}^{(k+1)} = \text{CoreLayers}\!\left(\mathbf{x}^{(k)}, \mathbf{S}^{(k)}\right), \quad k = 0, \ldots, K-1
\end{equation}

During training, $K$ is sampled uniformly from $\{1, \ldots, K_{\max}\}$ for each batch. For hardware compilation (e.g., on Tenstorrent accelerators), the loop is always unrolled to $K_{\max}$, with iterations beyond the sampled $K$ masked to identity via a blending operation:
\begin{equation}
    \mathbf{x}^{(k+1)} = \alpha_k \cdot \mathbf{x}_{\text{new}}^{(k)} + (1 - \alpha_k) \cdot \mathbf{x}^{(k)}
\end{equation}
where $\alpha_k = 1$ if $k < K$ and $\alpha_k = 0$ otherwise. This keeps the computation graph static while allowing variable-depth training.

\paragraph{$1/\sqrt{K}$ residual scaling.}
When the core layers are shared across $K$ iterations, gradients from all iterations accumulate on the shared parameters. Without normalization, the accumulated gradient norm scales as $\sim\sigma\sqrt{K}$, causing instability at large $K$. To normalize this, the blend factor is scaled by $1/\sqrt{K}$:
\begin{equation}
    \alpha_k = \frac{\mathbb{1}[k < K]}{\sqrt{K}}
\end{equation}
This turns the full replacement ($\alpha_k = 1$) into a partial update, where each iteration contributes $1/\sqrt{K}$ of the new state. The total variance added to the signal after $K$ iterations is then independent of $K$, following the variance-preserving initialization principle used in DeepNorm \citep{wang2022deepnorm}. The scaling uses the actual $K$ sampled for each batch, so batches with different $K$ values produce gradients of comparable magnitude.

At inference time, $K$ can be chosen freely---this is the test-time compute scaling mechanism. Harder problems can be allocated more iterations, directly trading compute for accuracy.

% ──────────────────────────────────────────────────────────────────────────
\subsection{Hybrid Backbone}

The full model uses 13--14 layers total, with attention inserted at positions 5 and 10. The attention layers provide exact retrieval capability---a known weakness of pure SSMs \citep{arora2024theoretical}---while the retention layers provide efficient sequence processing with flat memory. The hybrid ratio ($\sim$85\% SSM, $\sim$15\% attention) follows the empirical finding that a minority of attention layers recovers most of the retrieval gap while preserving the throughput and memory advantages of SSMs.

% ══════════════════════════════════════════════════════════════════════════
\section{Comparison to Thinking Tokens}
\label{sec:cot}
% ══════════════════════════════════════════════════════════════════════════

The dominant approach to scaling reasoning in language models is chain-of-thought (CoT) prompting: the model generates a sequence of intermediate ``thinking tokens'' before producing its final answer. This approach, while effective, has several structural inefficiencies that the workspace-recurrent architecture addresses directly.

\begin{table}[h]
\centering
\small
\begin{tabularx}{\textwidth}{p{0.2\textwidth} p{0.38\textwidth} p{0.38\textwidth}}
\toprule
\textbf{Property} & \textbf{Thinking Tokens (CoT)} & \textbf{Workspace Recurrence (Jasper)} \\
\midrule
Reasoning medium & Token space (discrete, serial) & Latent space (continuous, parallel) \\
Memory growth & Linear in reasoning length (KV cache) & Flat (fixed workspace + SSM state) \\
Tokens generated & One per reasoning step & Zero---reasoning is internal \\
Legibility & Fully legible (readable text) & Partially legible (workspace probes) \\
Compute control & Number of generated tokens & Number of loop iterations $K$ \\
Training data & Requires CoT supervision or RL & Standard next-token prediction + deep supervision \\
Expressive range & Limited to verbalizable reasoning & Can capture non-verbal reasoning patterns \\
Deployment cost & High (long generation, large cache) & Low (short output, fixed memory) \\
\bottomrule
\end{tabularx}
\caption{Comparison of thinking tokens vs.\ workspace recurrence.}
\label{tab:cot-comparison}
\end{table}

The key distinction is that thinking tokens reason in \emph{token space}---the model must verbalize each intermediate step, serialize it, and re-process it on the next forward pass. This is slow (one forward pass per token), memory-hungry (the KV cache grows with each token), and constrains reasoning to what can be expressed in language. As \citet{geiping2025recurrent} note, ``a substantial amount of thought happens through complex, recurrent firing patterns in the brain, before the first word of an answer is uttered.''

Workspace recurrence reasons in \emph{latent space}---the workspace state is updated in-place through matrix operations, without generating any tokens. This is faster (one forward pass per iteration, not per token), memory-flat (the workspace has a fixed size regardless of reasoning depth), and can capture reasoning patterns that are not easily expressed in words. The trade-off is legibility: thinking tokens produce readable traces, while workspace reasoning is only partially legible through probes (see \Cref{sec:experiments}).

For deployment on constrained devices (web browsers, mobile phones), the memory advantage is decisive. A model that generates 1000 thinking tokens accumulates a KV cache of $\sim$1000 $\times$ $d_{\text{model}}$ $\times$ $n_{\text{layers}}$ floats. A model that iterates its workspace 100 times uses the same fixed workspace state ($m \times d$ floats) plus the SSM state ($d_{\text{state}} \times d$ floats), regardless of $K$.

% ══════════════════════════════════════════════════════════════════════════
\section{Experimental Design}
\label{sec:experiments}
% ══════════════════════════════════════════════════════════════════════════

% ──────────────────────────────────────────────────────────────────────────
\subsection{Ablation Ladder}

We train three parameter-matched model variants ($\sim$10M parameters each) that form an ablation ladder---each cell adds one component, so the marginal contribution of each is isolated:

\begin{table}[h]
\centering
\small
\begin{tabularx}{\textwidth}{c l X}
\toprule
\textbf{Cell} & \textbf{Architecture} & \textbf{What it tests} \\
\midrule
A & Hybrid (14 retention + 2 attention at positions 5, 10), no workspace & The no-workspace control; isolates the workspace effect \\
B & Hybrid + workspace (13 retention + 2 attention + 16-slot perceiver workspace) & Does an engineered workspace help without recurrence? \\
C & Full architecture (hybrid + workspace + layers 6--9 looped $K$ times) & The go/no-go cell; does recurrence + workspace beat the control? \\
\bottomrule
\end{tabularx}
\caption{The three-cell ablation ladder. All cells target $\sim$10M parameters ($d_{\text{model}}=384$, 13--14 layers). Cell B removes one retention layer to compensate for the $\sim$2M workspace parameters, keeping cells parameter-matched.}
\label{tab:cells}
\end{table}

\paragraph{Cell history.}
The original design used a four-cell grid (A=pure Mamba-2, B=hybrid, C=hybrid+workspace, D=hybrid+workspace+recurrent) plus a learning-rate control (E). The pure-Mamba-2 cells (old A, old B) were deprecated after they plateaued at 0\% Task 1 accuracy, and the remaining cells were renamed: old E$\to$A, old C$\to$B, old D$\to$C. The SSM layer was also changed from Mamba-2's selective scan to a retention (decayed linear attention) layer, as described in \Cref{sec:coincidence}. The comparison logic is now: \textbf{B vs A} isolates the workspace effect (same LR=$2 \times 10^{-4}$); \textbf{C vs B} isolates the recurrent core; \textbf{C vs A} is the full R1 test.

% ──────────────────────────────────────────────────────────────────────────
\subsection{Training Configuration}
\label{sec:training-config}

All cells use AdamW ($\beta_1=0.9, \beta_2=0.95, \epsilon=10^{-8}$) with weight decay $0.1$ and gradient clipping at norm $1.0$. All cells use peak lr=$2 \times 10^{-4}$ with a linear warmup over 200 steps, followed by a constant peak. The training budget is 10{,}000 steps per cell. Training runs on a Tenstorrent Quietbox 2 (4$\times$ Blackhole chips, bfloat16-native) with all three cells in parallel on separate devices.

\paragraph{Batch configuration.}
Each step uses micro-batch size 8 with 48 gradient accumulation steps, giving an effective batch size of 384 sequences ($\sim$49K tokens per step). Fresh data is sampled every batch---with generated data there is no train set to overfit.

\paragraph{Early stopping.}
An EMA-smoothed loss plateau detector stops training if the loss does not improve by a relative $0.1\%$ over 1{,}000 steps (10\% of the training budget). The EMA smoothing factor is $\beta=0.99$, giving a half-life of $\sim$69 steps ($\sim$3{,}300 micro-batches), which is well above the per-step noise floor at this batch size.

\paragraph{Per-parameter LR groups.}
The workspace parameters ($\sim$2M of $\sim$10M total) receive $0.25\times$ the base learning rate ($5 \times 10^{-5}$ vs $2 \times 10^{-4}$ for the backbone). This is motivated by the observation that the workspace creates a sharper loss landscape: at equal LR, workspace gradients dominate and cause instability. The backbone LR is kept at $2 \times 10^{-4}$ to match Cell A (the control), preserving the clean B-vs-A comparison. Gate parameters (read/write gates) receive $10\times$ the workspace LR ($2 \times 10^{-3}$) so they open fast enough to contribute within the training budget.

\paragraph{Four stabilization techniques.}
The workspace and recurrent core introduce four sources of gradient instability that are not present in the baseline hybrid model. We address each with a targeted normalization:

\begin{enumerate}[nosep]
    \item \textbf{Per-parameter LR groups} (above): workspace parameters receive $0.25\times$ the base LR, preventing the workspace from dominating the gradient.
    \item \textbf{$1/\sqrt{K}$ residual scaling} (\Cref{sec:recurrent-core}): scales the recurrent core's blend factor by $1/\sqrt{K}$, normalizing gradient accumulation across $K$ iterations so that the total gradient norm is independent of $K$.
    \item \textbf{Slot parameter normalization} (\Cref{sec:workspace}): constrains the learned slot parameters to unit RMS after each optimizer step, breaking the positive feedback loop where large slots $\to$ sharp attention $\to$ large gradients $\to$ larger slots.
    \item \textbf{Spectral normalization} (\Cref{sec:workspace}): caps the spectral norm of all eight workspace weight matrices at $C=5$ after each optimizer step, addressing the root cause of unbounded weight growth directly while preserving sufficient attention selectivity (max ratio $\sim$164:1).
\end{enumerate}

Each technique targets a distinct mechanism: LR groups address the workspace's sharper loss landscape, $1/\sqrt{K}$ scaling addresses recurrent gradient accumulation, slot normalization addresses unbounded slot parameter growth, and spectral normalization addresses unbounded weight matrix growth. All four are deterministic (one hyperparameter: the spectral norm bound $C=5$) and can be disabled independently for ablation.

\paragraph{Custom fused kernels.}
The Tenstorrent implementation uses three custom tt-metal kernels to reduce op dispatch overhead and eliminate intermediate DRAM writes: (1)~fused RoPE (full rotation in a single kernel pass, using a split-RoPE optimization that avoids sub-tile concat/slice when $d_{\text{head}}/2$ is not tile-aligned), (2)~fused scale+decay (scores $= \text{qk} \cdot \text{scale} \cdot D_{\text{decay}}$ in one pass), and (3)~fused gate backward (computing $\text{grad\_out\_flat}$ and $\text{grad\_g}$ in a single kernel pass). The retention layer runs at 2.9\,ms per forward+backward pass at $d_{\text{model}}=384$, $n_{\text{heads}}=4$, $T=128$, $B=8$.

\paragraph{Attention regularizers (implemented, measured harmful, disabled).}
To counter workspace degeneracy we implemented an attention entropy penalty and a slot diversity penalty (\texttt{ws\_entropy\_weight}, \texttt{ws\_diversity\_weight}). Both are \textbf{disabled}, at weight 0. Enabling them at 0.01 reduced write-attention entropy by an order of magnitude while \emph{halving} task accuracy relative to the same architecture without them, because entropy rewards peakedness rather than input-dependence and the cheapest low-entropy solution is a fixed routing that ignores the input (\Cref{sec:results}). The code paths are retained for reproducibility of that negative result. Selectivity is now measured, not optimized: see the input-dependence diagnostic below.

\paragraph{Selectivity and decoding diagnostics.}
Two measurements replace the discarded entropy objective. \textbf{Input-dependence}: probe the model with $N=16$ distinct inputs of identical token length and compute the mean total variation distance between each input's attention distribution and the across-input mean, holding the position index fixed, plus the fraction of positions whose top-1 target is identical across all probes. A fixed routing scores TV $=0$ and 100\% agreement no matter how peaked it is, so this cannot be gamed by sharpening. Reported per path (read and write) because a functioning read path can otherwise mask a degenerate write path. \textbf{Teacher-forced answer accuracy}: evaluate answer tokens given the gold prefix, alongside free generation, to separate reasoning failure from exposure bias.

\paragraph{Architectural stability: zero-init gates and slot decay.}
The external constraints and regularizers address symptoms of instability but not the root cause: the workspace's feedback loop is amplifying rather than contractive. Two architectural changes make the workspace stable by construction. \textbf{Zero-initialized gates}: gate parameters are initialized to $-2$ ($\sigma(-2) \approx 0.12$), so the workspace starts with a small contribution and the feedback loop is mostly inactive. A gate-specific LR group (10$\times$ workspace LR) ensures the gates open within the training budget. \textbf{Slot decay}: the slot update includes a learned decay factor $\gamma$ (initialized at 1.0): $\text{slots} = \text{norm}(\gamma \cdot \text{slots} + g_{\text{read}} \cdot \text{read\_out})$. The decay provides a restoring force---old information fades unless actively reinforced---making the feedback loop contractive. Config: \texttt{gate\_init} (default $-2$), \texttt{slot\_decay\_init} (default $1.0$).

% ──────────────────────────────────────────────────────────────────────────
\subsection{Synthetic Tasks}

All tasks are character-level, generated on-the-fly (infinite data, no overfitting), and automatically verifiable. Each has a \textbf{depth knob} that controls how many reasoning steps a correct answer requires.

\paragraph{Task 1: Chained assignment arithmetic (multi-hop composition).}
Variables are defined in terms of previous variables, all arithmetic mod 97:
\begin{center}
\texttt{a=7;b=a*3+2;c=b-a;d=c*2;?d;}
\end{center}
The model must follow a chain of $k$ dependencies. Distractor variables are inserted at random positions. This is the core multi-hop reasoning task.

\paragraph{Task 2: Permutation tracking (state-tracking).}
$n$ items undergo $k$ swap operations; the model must report one item's final position:
\begin{center}
\texttt{n=6;2,5;0,3;5,1;?3;}
\end{center}
State-tracking is a documented weakness of linear SSMs \citep{arora2024theoretical}; this task tests whether the workspace compensates for the architecture's known gap.

\paragraph{Task 3: Single-hop recall (control).}
Shallow lookup with distractors:
\begin{center}
\texttt{a=5;b=12;c=8;d=3;e=15;?c;}
\end{center}
This is the control for selective ablation: the \jspace{} signature requires that killing the workspace leaves this intact while Tasks 1--2 collapse.

Training uses depths 2--8; evaluation extends to 2--16, measuring extrapolation beyond the training distribution.

% ──────────────────────────────────────────────────────────────────────────
\subsection{Pre-Registered Results}

Four measurements constitute proof, in ascending order of importance:

\begin{description}[style=multiline,labelwidth=3cm,leftmargin=3.2cm]
    \item[R1 --- Capability:] Cell C beats Cell A by $\geq$10 accuracy points on Tasks 1--2 at depths 10--16 (beyond training range).
    \item[R2 --- Compute scaling:] Accuracy on deep problems increases monotonically with $K$; harder depths need higher $K$ to saturate.
    \item[R3 --- Workspace decodability:] Linear probes on workspace slots can decode intermediate variable values, with decodability rising across loop iterations. This is the J-lens analogue.
    \item[R4 --- Selective ablation:] Replacing workspace slots with their training-set mean causes Tasks 1--2 to collapse ($\geq$30pt drop) while Task 3 barely moves ($\leq$5pt drop)---mirroring Anthropic's \jspace{} finding.
\end{description}

R3 and R4 are what make this a \emph{\jspace{} proof} rather than just another architecture ablation. R4 in particular tests the causal claim: the workspace is not merely correlated with reasoning performance, it is \emph{causally responsible} for it.

% ══════════════════════════════════════════════════════════════════════════
\section{Preliminary Results}
\label{sec:results}
% ══════════════════════════════════════════════════════════════════════════

The current training run (August 2026) uses the retention-based Mamba-3 layer on Tenstorrent Blackhole hardware with the three-cell grid (A, B, C) described above. All three cells are training in parallel on separate devices with a 10{,}000-step budget and early stopping. The following sections first summarize the historical results from the prior Mamba-2 runs (which motivated the architectural changes), then report the current run's early progress.

\paragraph{Current run status.}
As of this writing, all three cells are in the warmup phase (steps 50--100 of 200 warmup steps). Loss is dropping steadily: Cell A at 1.57 (step 100), Cell B at 1.56 (step 100), Cell C at 2.27 (step 50). Cell C is $\sim$2.7$\times$ slower per step (11.3\,s vs 4.1\,s) due to the recurrent core's $K_{\max}=6$ iterations. Estimated total training time is $\sim$12 hours for A/B and $\sim$31 hours for C. No evaluation results are available yet---the first eval is at step 500.

\paragraph{Historical context: Mamba-2 runs (deprecated).}
The following observations are from the prior Mamba-2 runs (old cells B, E, C, D), which have been deprecated. They are retained because they motivated the architectural changes (zero-init gates, slot decay, spectral normalization) and the methodological finding about entropy regularizers. The current Mamba-3 run starts fresh with these fixes incorporated from the beginning.

\paragraph{The no-workspace control plateaued on Task 1 (historical).}
Cell B (old hybrid baseline, lr=$6 \times 10^{-4}$) and Cell E (same architecture, lr=$2 \times 10^{-4}$) both plateaued at 0--5\% Task 1 accuracy---chance level---despite reaching 60--62\% overall (80\% Task 2, 100\% Task 3). The 3$\times$ learning rate difference did not change the Task 1 plateau, confirming it is an architectural limitation, not a learning-rate artifact. This is the benchmark the workspace cells must beat. In the current grid, Cell A serves as this control.

\paragraph{Workspace gradient instability and the four stabilization techniques (historical).}
Initial Cell C training (without stabilization, Mamba-2) collapsed at step 200 when the learning rate reached its peak: the gradient norm spiked from $\sim$1.3 to $\sim$13.4, and the model's output degenerated to a single token regardless of input. Cell E, which has no workspace, passed the same step without issue (grad norm 2.1), confirming the instability is workspace-specific.

We identified four distinct instability mechanisms and addressed each with a targeted normalization (\Cref{sec:training-config}):
\begin{enumerate}[nosep]
    \item \textbf{Workspace sharper loss landscape:} The workspace parameters dominate the gradient at equal LR. Per-parameter LR groups (workspace at $0.25\times$) delayed the collapse from step 200 to step 250 but did not prevent it.
    \item \textbf{Recurrent gradient accumulation:} Cell D's gradient norm was $\sim$20--22 at early steps (vs $\sim$5 for C/E), consistent with $\sqrt{K}$ accumulation. The $1/\sqrt{K}$ residual scaling reduced D's gradient norm to $\sim$9.7---roughly half, matching the theoretical prediction.
    \item \textbf{Unbounded slot growth:} The slot parameters grow via a positive feedback loop (large slots $\to$ sharp attention $\to$ large gradients $\to$ larger slots). Slot parameter normalization constrains slots to unit RMS after each optimizer step. With this fix, Cell C passed step 200 (the original collapse point) with grad norm 1.0, but collapsed at step 300 with grad norm 7.3---the slot vectors were constrained, but the attention weight matrices continued to grow.
    \item \textbf{Unbounded attention weight growth:} Even with slot magnitudes constrained, the $W_Q$ and $W_K$ projections can grow, producing sharp softmax distributions. We initially tried attention logit clamping ($C=5$) to cap the resulting softmax sharpness, but this was a symptom-level fix---the weights continued to grow, and Cell D still showed escalating gradient spikes (10.8 $\to$ 13.7). Spectral normalization at bound $C=1$ (the GAN default) eliminated the early spikes and allowed both cells to pass their previous collapse points (Cell C at step 300, Cell D at step 250 with grad norm 4.2, vs 15.2 in the 3-fix run). However, the collapse re-emerged at peak LR (step 400 for C with grad norm 19.7, step 300 for D with grad norm 61.1). Analysis revealed a \emph{capacity mismatch}: spectral norm 1 limits the attention ratio to $e^{2/\sqrt{d_h}} \approx 1.2{:}1$ (nearly uniform), so the workspace cannot provide the selective attention the backbone expects at peak LR. The backbone adapts fast (LR $2 \times 10^{-4}$), develops expectations that the workspace will provide selective attention, but the workspace is trapped at spectral norm 1 and cannot deliver---the resulting tug-of-war at the constraint boundary produces the gradient spike. The fix is to cap at $C=5$ instead of 1, allowing attention ratios up to $\sim$164:1---enough for selective attention, but still bounded. The cap (rather than hard normalization to $C$) lets weights grow naturally up to the bound, preventing disruption when resuming from checkpoints saved with smaller spectral norms.
\end{enumerate}
The first three stabilizations delayed the collapse (from step 200 to step 300 for Cell C) but did not prevent it. Spectral normalization at $C=1$ passed the old collapse points but re-collapsed at peak LR due to the capacity mismatch. The current approach---spectral norm capped at $C=5$---addresses the root cause while preserving sufficient attention capacity. These fixes are incorporated from the start in the current Mamba-3 run.

\paragraph{Workspace degeneracy and attention regularizers (historical).}
Even with stable training (no gradient collapse), diagnostic inspection of the Mamba-2 workspace at step 100 revealed a \emph{degeneracy} problem: the workspace was acting as a global average pool rather than a selective memory. Read attention entropy was 3.79 (out of $\ln 128 \approx 4.85$ for uniform), meaning slots attended nearly uniformly across all positions. Write attention entropy was 2.75 (out of $\ln 16 \approx 2.77$ for uniform), meaning every position wrote equally to every slot. The slot states did not change across recurrent core iterations. The workspace contained a blurred average of the entire sequence, not specific variable values.

This degeneracy arises because the cross-entropy loss provides only a diffuse gradient signal for workspace selectivity: the loss is several steps removed from the workspace's internal dynamics, and if the workspace isn't storing useful intermediates, the gradient for ``you should have stored $b$'' is weak. The model gets the same loss whether the workspace averages or selects, so there is no pressure to become selective.

We addressed this with two task-agnostic regularizers that can remain active during language training (unlike an auxiliary loss that would require known intermediate values):
\begin{enumerate}[nosep]
    \item \textbf{Attention entropy penalty} ($\lambda_{\text{ent}} = 0.01$): minimizes the mean entropy of read and write attention distributions, pushing the workspace toward selective (low-entropy) attention. This directly targets the observed uniformity without specifying \emph{what} to attend to.
    \item \textbf{Slot diversity penalty} ($\lambda_{\text{div}} = 0.01$): minimizes the mean cosine similarity between pairs of slots, pushing the 16 slots to carry different information rather than collapsing to identical averaged representations.
\end{enumerate}
Both regularizers are computed via host-side PyTorch autograd on the cached attention and slot tensors, with gradients injected into the workspace backward pass at the softmax and slot-state boundaries. They contribute $\sim$5\% of the total loss signal---a gentle nudge, not a dominant force.

\paragraph{The entropy regularizer is actively harmful: entropy is the wrong objective (historical).}
With the entropy and diversity regularizers active (weight 0.01 each), the Mamba-2 Cell C's write attention entropy fell from 2.75 to 0.224 (out of $\ln 16 \approx 2.77$ for uniform) and the answer-token loss reached 0.55---the lowest we have recorded, against 0.94 for the plateaued control. We initially read this as the workspace learning selective memory. Direct evaluation of that checkpoint shows the opposite.

Evaluated at 120 examples, the step-400 checkpoint scores 12.5\% overall and 5\% on Task 1, against 62\% overall for the no-workspace control (Cell E) and 33\% overall / 85\% Task 2 for the \emph{same architecture without regularizers}. Generated answers are near-constant regardless of input (`19', `13', `18', `18', `19', $\ldots$). The regularizers destroyed the only metric on which the workspace had ever exceeded the control.

The cause is a measurement error on our part, and it is the substantive methodological finding of this section. \textbf{Entropy measures peakedness, not input-dependence.} The cheapest way to minimize attention entropy is a \emph{fixed} one-hot routing that ignores the input entirely: such a pattern has near-zero entropy and carries zero information about the input. Our regularizer therefore selected precisely for the degenerate solution it was intended to prevent. Probing the step-400 checkpoint with 16 distinct same-length inputs and measuring the mean total variation distance of each attention pattern from the across-input mean confirms this directly:
\begin{center}
\begin{tabular}{lccc}
\toprule
\textbf{Path} & \textbf{Entropy} & \textbf{Input-dep.\ (TV)} & \textbf{Top-1 routing fixed} \\
\midrule
Write & 0.08--0.33 (very sharp) & 0.016 & \textbf{98.3\%} \\
Read  & 1.86--3.52 & 0.124 & 7.8--28.1\% \\
\bottomrule
\end{tabular}
\end{center}
For 98.3\% of positions, the write path routes to the \emph{same slot for every input}. Information enters the slots independently of content: the workspace had become a fixed positional bucket, not a content-addressed memory. The read path retained partial input-dependence (TV 0.124), which is why an aggregate verdict initially masked the failure---we now report input-dependence per path.

Two corrections follow, both adopted. First, selectivity must be measured as input-dependence (TV distance across probe inputs, plus top-1 routing agreement), never as entropy; the entropy and diversity regularizers are disabled. Second, the answer-token loss is a poor headline metric here: the label mask spans the answer digits \emph{and} the terminating EOS, so for a two-digit answer roughly one third of the supervised positions are trivially predictable, deflating the reported loss. Loss 0.55 coexisting with 12.5\% accuracy is largely explained by this dilution.

\paragraph{The Task 1 failure is a reasoning failure, not exposure bias (historical).}
Free-running generation is autoregressive, so a single wrong digit corrupts every subsequent step; a low generation score therefore conflates inability to reason with inability to recover from one's own errors. To separate these we now report \emph{teacher-forced} answer accuracy alongside generation, feeding the gold prefix at every answer position. On the step-400 checkpoint:
\begin{center}
\begin{tabular}{lccc}
\toprule
\textbf{Task} & \textbf{Generation} & \textbf{Teacher-forced (exact)} & \textbf{Teacher-forced (per token)} \\
\midrule
1 (chain arithmetic) & 5.0\% & 2.5\% & \textbf{22.4\%} \\
2 (permutation)      & 27.5\% & 47.5\% & 47.6\% \\
3 (single-hop recall) & 5.0\% & 12.5\% & 36.4\% \\
\bottomrule
\end{tabular}
\end{center}
Task 1 teacher-forced per-token accuracy is 22.4\% against a 10\% chance baseline over ten digits, and exact-match is 2.5\%. Even handed a perfect prefix, the model cannot produce the answer. This rules out exposure bias as the explanation for the Task 1 plateau and localizes the failure to the computation itself. (Task 2 does show a real exposure-bias gap, 47.5\% teacher-forced versus 27.5\% generated, so the diagnostic discriminates as intended rather than reporting a uniform verdict.)

\paragraph{Architectural stability: the feedback loop problem (historical).}
Analysis of the collapse mechanism reveals a fundamental architectural issue. The workspace has a \emph{feedback loop with gain}: it reads from $x$, updates slots, then writes from slots back to $x$. Each pass through this loop can amplify---sharp attention $\to$ large updates $\to$ changed representations $\to$ sharper attention. The external constraints (spectral norm, slot normalization, entropy regularization) cap individual components but do not make the loop itself contractive. This is structurally different from architectures that are stable by construction:
\begin{itemize}[nosep]
    \item \textbf{LSTMs}: the forget gate makes the cell state update contractive---the state can only change by a bounded amount per step.
    \item \textbf{ResNets/LoRA}: the residual branch is zero-initialized, so the module starts as identity and gradually learns to contribute.
    \item \textbf{Transformers}: each layer has residual + layer norm, and the attention output is bounded by the norm-constrained value vectors.
\end{itemize}

Our workspace has residual connections and norms, but two structural features prevent it from being naturally stable:

\textbf{(1) Gates initialized at 0.5.} The read and write gates use $\sigma(0) = 0.5$, so the workspace starts by modifying the state significantly rather than starting as identity. The feedback loop is active from step 0, before the model has learned to use it safely. This is the opposite of the zero-initialization principle that stabilizes additive modules.

\textbf{(2) No forget/decay on slots.} The slot update is $\text{slots} = \text{norm}(\text{slots} + g_{\text{read}} \cdot \text{read\_out})$---purely additive (plus normalization). There is no mechanism to decay or forget stored information. Once information enters a slot, it persists, and sharp attention keeps adding more. The feedback loop has no restoring force: if attention gets too sharp and overwrites a slot, nothing pulls it back.

\paragraph{Architectural fix: zero-init gates and slot decay.}
We address both issues with two structural changes that create a natural basin of attraction---a ``tennis ball in a bucket'' rather than a knife edge:

\textbf{(1) Zero-initialized gates.} Initialize the gate parameters to $-2$, giving $\sigma(-2) \approx 0.12$. The workspace starts with a small but non-trivial contribution: the feedback loop is mostly inactive. The gates gradually open as the model learns that the workspace is useful. A gate-specific LR group (10$\times$ the workspace LR) ensures the gates reach $\sigma(0) = 0.5$ within $\sim$1,000 steps, fitting the 3,000-step training budget. This is the same principle as zero-initialization in LoRA, tuned for a shorter training horizon.

\textbf{(2) Slot decay.} Change the slot update to $\text{slots} = \text{norm}(\gamma \cdot \text{slots} + g_{\text{read}} \cdot \text{read\_out})$ where $\gamma$ is a learned scalar initialized at 1.0. The decay makes the slot update contractive: old information naturally fades unless the read gate actively reinforces it. The feedback loop now has a restoring force---if attention gets too sharp and overwrites a slot, the decay pulls it back toward the learned slot embedding. The model must learn to overcome the decay, which means it only uses the workspace when the gradient signal is strong enough to justify it.

Together, these make the workspace start stable (near-zero gates) and stay stable (decaying slots). The external constraints (spectral norm, slot normalization, regularizers) become less critical because the architecture itself is contractive rather than amplifying.

\paragraph{Backward pass verification.}
The workspace backward pass has been verified against PyTorch autograd on CPU. All gradients---input gradients, weight gradients (8 matrices), slot gradients, norm weights, and gate scalars---match to within float32 precision ($<10^{-4}$ relative error). The gradient explosion was confirmed to be a training dynamics issue, not a backward pass bug.

\paragraph{Summary table (historical, Mamba-2 runs).}
\Cref{tab:eval-results} shows the evaluation results from the deprecated Mamba-2 runs, collected by the automated eval loop. These results motivated the architectural changes (zero-init gates, slot decay, $C=5$ spectral norm) and the methodological finding about entropy regularizers. The current Mamba-3 run incorporates all these fixes from the start; its results will be reported once training completes.

\begin{table}[ht]
\centering
\small
\begin{tabular}{c c c c c c}
\toprule
\textbf{Cell} & \textbf{Step} & \textbf{Overall} & \textbf{Task 1} & \textbf{Task 2} & \textbf{Task 3} \\
\midrule
A & 200 & 30\% & 0\% & 80\% & 10\% \\
A & 300 & 35\% & 0\% & 75\% & 30\% \\
B & 100 & 22\% & 0\% & 60\% & 5\% \\
B & 200 & 30\% & 5\% & 80\% & 5\% \\
B & 300 & 57\% & 0\% & 75\% & 95\% \\
B & 400 & 60\% & 0\% & 80\% & 100\% \\
E & 100 & 22\% & 0\% & 60\% & 5\% \\
E & 200 & 32\% & 0\% & 75\% & 20\% \\
E & 300 & 28\% & 0\% & 70\% & 15\% \\
E & 400 & 38\% & 0\% & 75\% & 40\% \\
E & 500 & 40\% & 0\% & 75\% & 45\% \\
E & 600 & 62\% & 5\% & 80\% & 100\% \\
\midrule
\multicolumn{6}{l}{\textit{Old runs (no stabilization, collapsed at step 200):}} \\
C (old) & 100 & 23\% & 0\% & 60\% & 10\% \\
C (old) & 200 & 7\% & 5\% & 15\% & 0\% \\
D (old) & 100 & 5\% & 0\% & 15\% & 0\% \\
D (old) & 200 & 22\% & 0\% & 60\% & 5\% \\
\midrule
\multicolumn{6}{l}{\textit{With LR groups + slot norm + $1/\sqrt{K}$ (collapsed at step 300):}} \\
C (3-fix) & 300 & 25\% & 0\% & 70\% & 5\% \\
C (3-fix) & 400 & 23\% & 0\% & 65\% & 5\% \\
C (3-fix) & 500 & 20\% & 0\% & 60\% & 0\% \\
C (3-fix) & 600 & 15\% & 0\% & 45\% & 0\% \\
D (3-fix) & 200 & 22\% & 0\% & 60\% & 5\% \\
\midrule
\multicolumn{6}{l}{\textit{With LR groups + slot norm + $1/\sqrt{K}$ + spectral norm $C=1$ (passed old collapse, re-collapsed at peak LR):}} \\
C ($C{=}1$) & 300 & 33\% & 0\% & 85\% & 15\% \\
D ($C{=}1$) & \multicolumn{5}{l}{re-collapsed at step 300 (grad norm 61)} \\
\midrule
\multicolumn{6}{l}{\textit{$C=5$ + entropy/diversity regularizers (harmful; 30 examples/task):}} \\
C (reg) & 400 & 12\% & 5\% & 28\% & 5\% \\
C (reg) & \multicolumn{5}{l}{answer-loss 0.55 (lowest recorded) yet worst eval; write routing 98.3\% input-independent} \\
\midrule
\multicolumn{6}{l}{\textit{Architectural fix (zero-init gates + slot decay), regularizers off, in progress:}} \\
C (arch) & \multicolumn{5}{l}{training, 3{,}000-step budget} \\
D (arch) & \multicolumn{5}{l}{training, 3{,}000-step budget} \\
\bottomrule
\end{tabular}
\caption{Evaluation results from the deprecated Mamba-2 runs (5 examples per task per depth unless noted, depths 2/4/6/8). Cell names use the old naming (A=pure Mamba-2, B=hybrid, C=hybrid+workspace, D=hybrid+workspace+recurrent, E=LR control). These results are historical---the current Mamba-3 run uses the renamed three-cell grid (A/B/C) with all stabilizations applied from the start. \textbf{No Mamba-2 configuration ever exceeded noise on Task 1}, with or without a workspace. The workspace's entire measured advantage over the control, at its historical best, was 85\% versus 80\% on Task 2---one task, five points, one seed.}
\label{tab:eval-results}
\end{table}

% ══════════════════════════════════════════════════════════════════════════
\section{Discussion}
\label{sec:discussion}
% ══════════════════════════════════════════════════════════════════════════

% ──────────────────────────────────────────────────────────────────────────
\subsection{Why This Might Work}

The \jspace{} finding tells us that reasoning in large models is carried by a small, privileged subspace. The TRM finding tells us that iteration through a small network can substitute for parameter count. The SSM coincidence tells us that the SSM's compressed state already has the functional shape of the \jspace{}. Jasper combines all three: the workspace provides the privileged subspace, the recurrent core provides the iteration, and the hybrid SSM backbone provides the substrate.

If this works, the implication is that we can build reasoning models that are:
\begin{itemize}[nosep]
    \item \textbf{Small enough for the browser:} $\sim$1B parameters, flat memory in both sequence length and loop count.
    \item \textbf{Controllable at inference:} $K$ can be adjusted per problem, trading compute for accuracy.
    \item \textbf{Interpretable:} workspace probes (R3) provide a window into the model's reasoning, unlike opaque latent reasoning in recurrent transformers.
    \item \textbf{Verifiable:} selective ablation (R4) provides a causal test of whether the workspace is doing the reasoning, not just correlated with it.
\end{itemize}

% ──────────────────────────────────────────────────────────────────────────
\subsection{Why It Might Not Work}

Synthetic-task wins at 30M parameters have a real history of not transferring to language at scale. The tasks we use (arithmetic chains, permutation tracking) are structurally simple and may not capture the full complexity of language reasoning. The workspace may help on these tasks without generalizing to the kind of reasoning that matters for real applications.

The gradient instability consumed most of our effort, and we should be clear that little of what resolved it was novel. Stacking $K$ residual updates accumulates gradient as $\sqrt{K}$, so $1/\sqrt{K}$ scaling is arithmetic; a parameter inside a positive feedback loop needs its magnitude constrained, so slot normalization is routine; unbounded attention logits call for spectral normalization, standard since \citet{miyato2018spectral}; an additive module should begin at identity, standard since LoRA; and a recurrent state with no forget term will not stay bounded, which LSTMs settled in 1997. The capacity mismatch at $C=1$ is likewise a closed-form consequence---bounding the spectral norm at $C$ caps attention logits at $C^2/\sqrt{d_h}$, giving a ratio of $\approx 1.2{:}1$ at $C=1$, which is uniform to within noise---and could have been computed before the run rather than diagnosed after it. These properties belong in the initial design, and we present them in \Cref{sec:training-config} as design requirements derived from the architecture rather than as findings. The cost of not deriving them up front was roughly eight restarts.

What was not obvious, and what we got wrong, is the measurement. We adopted attention entropy as a proxy for selectivity and optimized it directly; because entropy scores peakedness rather than input-dependence, the objective was minimized by a routing that ignores the input, and the resulting checkpoint posted our best loss and our worst accuracy simultaneously. The general lesson---that a proxy admitting a degenerate optimum will find it, and that ``sharp'' must be distinguished from ``informative''---transfers beyond this architecture, and it is the part of this work we would defend.

The more substantive question is whether the workspace provides enough \emph{computational leverage} to matter, and here we must be blunt: \textbf{in the Mamba-2 runs, there was no evidence that it did.} Across every Mamba-2 configuration we trained---unstabilized, three stabilizations, spectral norm $C=1$, $C=5$, and $C=5$ with regularizers---Task 1 accuracy never exceeded noise, for workspace and non-workspace cells alike. The workspace's entire measured advantage over the control at its historical best was 85\% versus 80\% on Task 2: one task, five points, a single seed, no error bars. Teacher-forcing ruled out decoding as the culprit (22.4\% per-token against a 10\% baseline), so the failure was in the computation. The current Mamba-3 run, with all stabilizations applied from the start and a 10{,}000-step budget (vs.\ the 3{,}000-step Mamba-2 budget), will provide the first clean test of whether the architectural fixes change this picture.

We also have to be honest that our most encouraging number was an artifact. The loss of 0.55 that appeared to show the workspace learning was produced by a checkpoint that evaluates at 12.5\%, and was driven by a regularizer optimizing a metric (entropy) that does not measure the property we cared about (input-dependence). That is a cautionary result about proxy objectives, and it is the strongest finding in this section.

Two possibilities remain open. The first is that our experimental design tests the wrong thing: Task 1 requires iteration \emph{and} memory, and Cell B supplies memory without iteration, so B-versus-A may be architecturally incapable of showing a benefit regardless of workspace quality. On this reading the informative comparison is C-versus-B, and we have under-prioritized C because it is slower and less stable. The second is that the workspace as designed is simply not the right mechanism. It requires eight weight matrices, two gates, a decay term, two normalizations, and three learning-rate groups to train at all; that quantity of scaffolding is itself evidence against the design, since the architectures that won in practice---transformers included---won by being simple and trainable rather than by being optimal. Accordingly we have set an advance stopping rule: if Cell C shows Task 1 $\leq 10\%$ by step 2{,}000 with gates verifiably open and attention verifiably input-dependent, we report the negative result rather than continue to search for one more fix.

Finally, the SSM coincidence may be more metaphorical than functional. The SSM state and the \jspace{} have similar shapes, but they operate at different levels of abstraction---the SSM state is a low-level sequence summary, while the \jspace{} is a high-level conceptual workspace. Whether the workspace module can bridge this gap is an empirical question.

% ──────────────────────────────────────────────────────────────────────────
\subsection{Decision Framework}

The proof-of-concept is designed to produce a clear go/no-go decision before any significant cloud spend:

\begin{itemize}[nosep]
    \item \textbf{Go:} R1 and R2 pass, and at least one of R3/R4 shows the workspace doing real causal work $\to$ fund a 300M language-scale ablation ($\sim$\$3--5K).
    \item \textbf{Pivot:} R1 fails but Cell A's efficiency story stands $\to$ keep the hybrid backbone, drop the workspace, proceed on hybrid distillation.
    \item \textbf{Kill:} Cell C $\leq$ Cell A everywhere $\to$ the novel architecture track is dead; fall back to sampled-attempts-plus-verifier on a standard model.
\end{itemize}

The entire proof-of-concept costs $\sim$\$20--30 of electricity and under a week of compute. This buys the right to spend \$5K intelligently, not the conclusion itself.

% ══════════════════════════════════════════════════════════════════════════
\section{Related Work}
\label{sec:related}
% ══════════════════════════════════════════════════════════════════════════

\paragraph{Global workspace in LLMs.}
\citet{anthropic2026jspace} discovered the \jspace{} in Claude Opus 4.6 using the Jacobian lens, showing that a small set of verbalizable representations functions as a global workspace \citep{baars2005global,dehaene2014conscious}. Our workspace module is an explicit engineering of this concept.

\paragraph{Iterative and recurrent reasoning.}
\citet{jolicoeur2025trm} showed that tiny recursive networks (7M parameters) can outperform much larger models on ARC-AGI through iterative refinement. \citet{geiping2025recurrent} scaled recurrent-depth transformers to 3.5B parameters, demonstrating test-time compute scaling in latent space. \citet{gehring2024relaxed} explored relaxed recursive transformers with shared parameters. Our recurrent core follows the TRM paradigm but uses a hybrid SSM/attention backbone and an explicit workspace state.

\paragraph{Hybrid SSM-attention models.}
\citet{dao2024mamba2} established the SSD framework connecting SSMs and attention. \citet{wang2025m1} showed that hybrid Mamba reasoning models can match transformer-based R1-distilled models on AIME/MATH while generating 3$\times$ faster. NVIDIA's Nemotron line has repeatedly shown hybrid models matching or beating same-size transformers on long reasoning traces. Our architecture follows the hybrid pattern but adds the workspace and recurrent core. The current implementation uses a retention (decayed linear attention) layer in place of the selective scan, following the RetNet formulation \citep{sun2024retention}, which preserves the SSM's fixed-size state property while being simpler to implement on the target hardware.

\paragraph{Test-time compute scaling.}
The broader literature on scaling test-time compute---from self-consistency \citep{wang2023selfconsistency} to process reward models \citep{lightman2023prm} to OpenAI's o1/o3 models---operates in token space. Our approach scales test-time compute in latent space, which is more memory-efficient but less legible.

% ══════════════════════════════════════════════════════════════════════════
\section{Conclusion}
% ══════════════════════════════════════════════════════════════════════════

We have presented Jasper, an architecture that combines three recent insights---the \jspace{} as a compact reasoning workspace, iterative refinement as a substitute for parameter count, and the SSM's compressed state as a native substrate for the workspace---into a single system. The architecture reasons in latent space without generating tokens, keeps memory flat in both sequence length and loop count, and provides interpretable and causally testable reasoning through workspace probes and selective ablation. The current implementation uses a retention (decayed linear attention) layer that preserves the SSM's key properties while being simpler and faster on the target hardware.

The proof-of-concept experiment is designed to produce a clear verdict at minimal cost: if the workspace causally carries multi-step reasoning (R4) and the recurrent core provides test-time compute scaling (R2), the architecture warrants scaling to language-scale experiments. If not, the hybrid Mamba-attention backbone still stands as an efficient baseline for in-browser deployment.

The broader claim is that reasoning in language models need not be tied to token generation. The brain reasons through recurrent firing patterns before articulating a thought; perhaps language models should too.

% ══════════════════════════════════════════════════════════════════════════
% References
% ══════════════════════════════════════════════════════════════════════════

\begin{thebibliography}{99}
\small

\bibitem[Anthropic(2026)]{anthropic2026jspace}
Wes~Gurnee and Jack~Lindsey et~al.
\newblock Verbalizable representations form a global workspace in language models.
\newblock \emph{arXiv preprint arXiv:2607.15495}, 2026.

\bibitem[Miyato et~al.(2018)]{miyato2018spectral}
Takeru Miyato, Toshiki Kataoka, Masanori Koyama, and Yuichi Yoshida.
\newblock Spectral normalization for generative adversarial networks.
\newblock In \emph{International Conference on Learning Representations (ICLR)}, 2018.

\bibitem[Jolicoeur-Martineau(2025)]{jolicoeur2025trm}
Alexandre Jolicoeur-Martineau.
\newblock Less is more: Recursive reasoning with tiny networks.
\newblock \emph{arXiv preprint arXiv:2510.04871}, 2025. [Samsung SAIT Montreal]

\bibitem[Geiping et~al.(2025)]{geiping2025recurrent}
Jonas Geiping, Sean~McLeish, Neel~Jain, John~Kirchenbauer, Siddharth~Singh, Brian~R.~Bartoldson, Bhavya~Kailkhura, Abhinav~Bhatele, and Tom~Goldstein.
\newblock Scaling up test-time compute with latent reasoning: A recurrent depth approach.
\newblock In \emph{Advances in Neural Information Processing Systems (NeurIPS)}, 2025.

\bibitem[Dao and Gu(2024)]{dao2024mamba2}
Tri~Dao and Albert~Gu.
\newblock Transformers are SSMs: Generalized models and efficient algorithms through structured state space duality.
\newblock In \emph{Proceedings of the 41st International Conference on Machine Learning (ICML)}, pages 10041--10071, 2024.

\bibitem[Baars(2005)]{baars2005global}
Bernard~J.~Baars.
\newblock Global workspace theory of consciousness: Toward a neuroscience of human experience.
\newblock \emph{Progress in Brain Research}, 150:45--53, 2005.

\bibitem[Dehaene(2014)]{dehaene2014conscious}
Stanislas~Dehaene.
\newblock \emph{Consciousness and the Brain: Deciphering How the Brain Codes Our Thoughts}.
\newblock Viking, 2014.

\bibitem[Wang et~al.(2025)]{wang2025m1}
Junxiong~Wang et~al.
\newblock M1: Towards scalable test-time compute with Mamba reasoning models.
\newblock \emph{arXiv preprint}, 2025.

\bibitem[Arora et~al.(2024)]{arora2024theoretical}
Simran~Arora, Sabhash~Rao, and Christopher~R\'e.
\newblock Theoretical limitations of self-attention and state space models in reasoning tasks.
\newblock \emph{arXiv preprint}, 2024.

\bibitem[Gehring et~al.(2024)]{gehring2024relaxed}
Jonas~Gehring, Kilian~Golde, and David~Grangier.
\newblock Relaxed recursive transformers.
\newblock \emph{arXiv preprint arXiv:2410.10672}, 2024. [Google DeepMind]

\bibitem[Wang et~al.(2023)]{wang2023selfconsistency}
Xuezhi~Wang et~al.
\newblock Self-consistency improves chain of thought reasoning in language models.
\newblock In \emph{International Conference on Learning Representations (ICLR)}, 2023.

\bibitem[Lightman et~al.(2023)]{lightman2023prm}
Hunter~Lightman et~al.
\newblock Let's verify step by step.
\newblock \emph{arXiv preprint arXiv:2305.20050}, 2023. [OpenAI]

\bibitem[Wang et~al.(2022)]{wang2022deepnorm}
Hongyu~Wang et~al.
\newblock DeepNet: Scaling transformers to 1,000 layers.
\newblock \emph{arXiv preprint arXiv:2203.00555}, 2022. [Microsoft]

\bibitem[Sun et~al.(2024)]{sun2024retention}
Yutao~Sun, Li~Dong, Shaohan~Huang, Shuming~Ma, Yuqing~Xia, Jilong~Xue, Zhangyin~Feng, and Furu~Wei.
\newblock Retentive network: A successor to transformer for large language models.
\newblock \emph{arXiv preprint arXiv:2307.08621}, 2024. [Microsoft]

\end{thebibliography}

\end{document}
